\chapter{The Hilbert Completion Door}

Every earlier chapter has named a door and declined to open it. The square root that
would turn an energy into a length; the completion that would close the normed activities
into a full geometry; the generic object the forcing conditions approach but never form ---
each was framed, marked, and left shut. This chapter walks up to that door, which is one
door under three names, and does the one thing a finite apparatus honestly can: it names,
precisely and in order, what lies beyond, and it states its final result as a conditional
that depends on exactly those named things. It does not open the door. It writes down what
opening it would require.

The chapter is, by design, the most cautious in the book. It claims no continuum
orthogonality, no real Hilbert completion, no standard theorem of the larger physics. What
it builds is an interface --- a precise list of the obligations the continuum theory would
have to meet --- and it shows that the finite three-rung object the book has carried can
stand at the entrance and present its credentials. The credentials are real: a
zero-detecting norm, a finite pairing, a countable ledger, the gauge residue. The
obligations are named: coercivity, completeness, density, a continuum pairing. The chapter
keeps the two lists strictly apart, the built from the assumed, so that no token can
be mistaken for a theorem.

The first section opens the door as an interface. It defines what a Sobolev-and-Hilbert
door must provide --- a normed space with zero detection, a pairing, and two obligations
carried as named tokens, coercivity and completeness --- and shows the finite object
standing at it, presenting its norm and its pairing and carrying the completeness token
without pretending to have discharged it. The door names the obligations; it does not meet
them.

The second section makes the skeleton that sits inside the door. The finite forcing ledger
of the previous chapter, with its countable coding, becomes a separable pre-Hilbert
skeleton: a countable dense set standing inside the completion. Here the chapter is exact
about which words are real and which are tokens. ``Countable'' is concrete --- every finite
condition carries an actual natural-number code. ``Dense'' and ``completion'' are tokens,
carried by the door and not discharged by the skeleton. The skeleton is genuinely
countable and named-dense, and the chapter says exactly that.

The third section names the certificate families that would live inside the completion. A
continuum certificate is a vector in the completed space together with a finite witness ---
a geometric program on one side, the generated gauge program on the other. The chapter is
candid that the vectors it can supply concretely are, for now, the zero vector: the finite
witnesses carry all the content, and a genuine nonzero dense embedding is named as the
analytic upgrade the door still awaits. The families are named with honest placeholders.

The fourth section lifts the finite structure of the boundary-radiation result to the
door, as three interfaces. Orthogonality becomes a statement about the continuum pairing;
the Clifford obstruction becomes a door-facing rule that orthogonality forbids an interior
mixed certificate; and the boundary trace becomes a functional that reads the flux and
agrees with the terminal residual on the generated representative. Each is the finite
result of the previous chapter, restated as an obligation the continuum theory would carry.

The fifth section states the chapter's --- and in a sense the book's --- conditional theorem.
Given the door-facing hypotheses, all of them named in the preceding sections --- generated
certificates, a continuum orthogonality, a compatible Clifford structure, a boundary trace
that agrees with the residual, and a Cauchy residual sequence in the completion norm --- the
generalized boundary-radiation result follows: no interior mixed certificate, flux equal to
the residual, and a convergent residue. The theorem is an implication, and the chapter says
so plainly. The finite object discharges the hypotheses in the only way it can, through the
zero-vector lift, and what remains genuinely analytic is left as the named antecedent of an
honest conditional.

By the end of the chapter the continuum is not conquered; it is described. The apparatus
stands at the door with its finite credentials in hand, names every obligation the far side
would impose, and states its last large result as a theorem conditional on exactly those
obligations. That is the honest shape of a finite theory at the edge of the continuum: not
a claim to have crossed, but a precise account of what crossing would cost.

\section{The Sobolev/Hilbert Door and the Finite Skeleton}

The previous chapter left the apparatus with a clean finite ledger and a
closed finite residue account, but not with a completed space. This section
opens Chapter 16 by naming the door that separates those two achievements. The
door is not a theorem that the finite apparatus has already proved. It is an
interface: a reader-facing list of the structures and hypotheses that would
allow the finite construction to be read inside Sobolev and Hilbert language.
The discipline of the chapter begins here, with the distinction between a
credential already built and an obligation only named.

On the built side, the apparatus brings the three-rung object, its
zero-detecting norm, its finite pairing, its countable ledger of finite
conditions, and the finite gauge residue left by the mixed pair. These are not
decorations. The norm is meant to detect the trivial vector: if its length
vanishes, the finite representative has collapsed to zero. The pairing is the
finite shadow of an inner product, symmetric in the relevant entries and
harmless on the zero vector. The ledger is countable because the preceding
chapter coded finite partial commitments by natural numbers. The residue is
finite because the vacuum/linear part has closed and the mixed-pair charge is
what remains. Those are the credentials presented at the threshold.

On the completion side, the door asks for more than the apparatus has built. It
asks for coercivity, so that the norm controls the object strongly enough to
support analytic comparison. It asks for completeness, so that Cauchy motion in
the norm has a place to land. It asks for density, so that the countable
skeleton is not merely countable but genuinely approximating. It asks for a
continuum pairing, so that the finite pairing can be read as the trace of a
completed Hilbert structure rather than only as a finite table. These are
obligations, not hidden conclusions. The section names them so that later
claims can say exactly where their analytic strength enters.

The door can therefore be summarized as a compact contract:
\[
  \begin{aligned}
  \text{door}=\big(&
    \text{zero-detecting norm},\ \text{finite pairing};\\
    &\text{coercivity},\ \text{completeness},\
      \text{density},\ \text{continuum pairing}
  \big).
  \end{aligned}
\]
The first two entries are supplied by the finite apparatus in the form needed
to stand at the door. The remaining entries are carried as named hypotheses for
the completion theory. This notation is intentionally plain. It is not trying
to smuggle a Hilbert space into the chapter by typography. It is recording the
minimum vocabulary required before the later sections may speak about a
completion without pretending that countable finite data alone created one.

The distinction is also a guard against a common false comfort. A finite norm
can detect zero without controlling every analytic mode that a completed space
would contain. A finite pairing can be symmetric without already being the
inner product of a complete Hilbert space. A countable ledger can enumerate
finite records without proving that those records approach every continuum
state. Each finite credential is valuable precisely because it is explicit and
checkable, but none of them silently acquires the missing analytic clauses. The
door is the place where those clauses are written in full view.

The finite skeleton is the part of the contract the apparatus can actually
place on the near side of the door. From Chapter 15 it inherits finite
conditions, extension, compatibility, merge, finite decision, extensional
equality, and natural-number coding. Those pieces give a countable family of
finite records. As a skeleton, the family may be placed inside the named
completion interface as the candidate dense set. But the word ``candidate''
matters. Countability is constructed. Density is a completion-side obligation.
The skeleton is a genuine finite object equipped with a name tag for the
analytic property it would need in order to become a separable pre-Hilbert
spine.

This is why the section does not yet teach the separable completion story. It
only prepares the ground for it. The next section can discuss how the
countable ledger sits as a pre-Hilbert skeleton precisely because this section
has already separated the earned finite ingredients from the named analytic
requirements. Without that separation, the phrase ``separable completion''
would arrive too quickly, as if a natural-number code for finite conditions
already produced a dense subset of a completed space. The chapter refuses that
shortcut.

Nor does the section claim the boundary-radiation theorem in continuum form.
The finite residue is present, and the finite pairing can remember it, but
continuum orthogonality, Sobolev completion, Hilbert completion, and boundary
trace all belong to later hypotheses. To stand at the Sobolev/Hilbert door is
therefore an honest achievement with a limited shape: the finite apparatus has
credentials in hand, the completion side has obligations named, and the two
lists are kept separate. The chapter will move forward only by respecting that
separation.

That limited shape is enough for the chapter's first step. The reader should
leave this section knowing exactly which side of the threshold each word
occupies. Norm and finite pairing belong to the apparatus. Countable coding
belongs to the apparatus. Coercivity, completeness, density, continuum pairing,
orthogonality, and boundary trace belong to the door-facing hypotheses. The
finite skeleton may be carried to the door, but the door remains a door: a
marked place where a completed theory must do work the finite apparatus has
not done for it.

\section{The Separable Pre-Hilbert Completion Skeleton}

The door contract of the previous section has two columns. On the near side are
finite credentials: a zero-detecting norm, a finite pairing, a countable ledger
of conditions, and the finite gauge residue. On the far side are obligations:
coercivity, completeness, density, continuum pairing, orthogonality, and
boundary trace. This section develops only the near-side object that can be
carried to that door. It is called a separable pre-Hilbert completion skeleton
because it has the formal shape that such a completion would need, not because
the completion has already been supplied.

The first earned word is ``countable.'' Chapter 15 did not merely assert that
the finite records can be listed. It built the listing out of the records
themselves. A finite condition is a finite set of commitments. Each commitment
has finitely described coordinates, a finite tolerance, and a finite value
description. A condition can therefore be read extensionally as a finite list of
finite data. Once proof-route decorations have been erased by extensional
equality, two conditions with the same commitments count as the same record for
the ledger. Finite lists of finite codes have natural-number codes, so the
ledger has the form
\[
  p\ \longmapsto\ \ulcorner p\urcorner \in \mathbb{N}.
\]
The notation is modest, but the achievement is real. The apparatus has not
guessed that its conditions are countable; it has made their countability part
of their receipt.

This is also why extensional equality from the previous chapter is not a
technical ornament. Without it, the same finite commitments could appear under
many descriptions, and the ledger would count performances rather than records.
The skeleton needs records. A countable pre-Hilbert candidate cannot be made
from the history of how a condition was found; it has to be made from the
condition's finite content. The grid image from the previous chapter is useful
here. Two half-filled grids that display the same letters in the same squares
are the same partial grid, even if one constructor reached it by working across
and another by working down. Likewise, two finite conditions with the same
commitments give the same skeleton record. The code belongs to the displayed
finite fact, not to the route by which the fact was obtained.

That countability is the beginning of separability, not the end of it.
Separability in the analytic sense requires a countable dense subset of a space.
The finite ledger supplies the countable family. It does not, by that fact
alone, supply the space in which density is measured, nor does it prove that its
records approximate every point of such a space. Density talks about distance
to arbitrary points of a completed object. Completion talks about limits of
Cauchy motion. Both notions require the far side of the door. The chapter
therefore keeps the sentence in two pieces: the countable family is built; its
density in a completed Hilbert or Sobolev object is a named obligation.

The word ``pre-Hilbert'' receives the same treatment. The apparatus has a
finite pairing and a norm-like zero detector. These give the finite ledger the
right kind of algebraic posture: records may be compared, paired, and assigned
length in the finite sense already available. That is enough to organize the
ledger as a pre-Hilbert skeleton. It is not enough to announce a Hilbert space.
A Hilbert space would require the completed normed object, the inner product
that survives completion, and the guarantee that limit processes land where
they are supposed to land. The skeleton has the finite grammar out of which
that language could be read. It does not contain the analytic closure of the
language.

This distinction is easiest to see by following a single finite record. A
condition \(p\) has a code \(\ulcorner p\urcorner\). It also has a finite
support, a finite list of commitments, and whatever finite pairing data the
door contract allows one to evaluate between it and another finite record. If
another condition \(q\) extends it, then \(q\) carries all of \(p\)'s
commitments and perhaps more. If two conditions are compatible, their finite
merge is still a condition. If a finite question must be decided, a finite
extension can carry the new commitment. These operations make the ledger
stable as a finite family. They also explain why the ledger is the right
candidate to present at the completion door: it has enough internal discipline
to be transported as a skeleton rather than as a heap of unrelated finite
marks.

The same operations explain the phrase ``completion skeleton.'' A completion
story needs approximants that can be compared and refined. Extension gives the
direction of refinement. Compatibility says when two partial approximants can
be held together. Merge gives the finite common refinement when there is no
conflict. Finite decision says that a finite question can be pushed into the
record by a finite extension. None of these operations says that a limit
exists. Taken together, however, they give the family the shape of approximants
for a possible limit theory. The skeleton is not a pile of samples; it is a
finite directed grammar that a completion hypothesis would know how to read.

But a candidate dense family is still a candidate. The apparatus can say:
here is the countable family; here are its finite extension and compatibility
operations; here is the finite pairing; here is the zero-detecting norm; here
is the residue that the finite gauge bookkeeping leaves. The apparatus cannot
say, from those facts alone: every point of the completed object is a norm
limit of this family. That last sentence is exactly where density enters. Nor
can it say: every Cauchy sequence of skeleton records converges inside an
already available space. That sentence is exactly where completeness enters.
The skeleton is therefore threshold-facing. It is prepared for density and
completion, but it does not impersonate them.

The finite gauge residue travels with this skeleton in the same constrained
way. Chapter 15 closed the vacuum/linear account and filed the mixed pair as
the remaining finite charge. The present section does not remake that
bookkeeping as an analytic theorem inside a completed space. It records how the
residue is carried by the finite family that stands at the door. Each skeleton
record may include the finite gauge information already earned: which
linear/vacuum entries have closed, where the mixed-pair residue remains, and
how that residue is paired with the geometric side of the construction. When
the ledger is coded, this gauge information is part of the finite receipt. When
the ledger is presented as a pre-Hilbert skeleton, the residue comes along as
finite mapped content, not as a new continuum claim.

That mapped content matters because the completion door is not only a place for
abstract vectors. It is where the earlier finite geometry and finite gauge
accounting must be held together without changing their status. If the
skeleton forgot the residue, the completion story would have cut away one of
the main facts that Chapter 15 delivered. If the skeleton promoted the residue
to a Sobolev boundary theorem, it would have crossed the door without paying
the analytic price. The section avoids both failures. It carries the residue as
finite data attached to a countable family and leaves the analytic reading of
that data to the hypotheses that would make the completion legitimate.

The point is subtle but important. A finite gauge residue can be transported in
two different ways. One way is bookkeeping: a skeleton record remembers that
the linear/vacuum part has closed and that the mixed pair remains. That is the
way available here. The other way is analytic interpretation: a completed
object would provide a normed location, a boundary trace, and an orthogonality
claim in which that residue could be read as a continuum phenomenon. The
section refuses to confuse the two. It permits the first transport because the
finite ledger already supports it. It withholds the second until the door
hypotheses have done their work.

The resulting object can be described in a compact way:
\[
  \mathcal{S}_{0}
    = \{\,\text{finite coded records with finite pairing and residue data}\,\}.
\]
The subscript \(0\) is a reminder of its status. This is the ground skeleton,
not the completed space. It is countable because its records are coded by
natural numbers. It is pre-Hilbert in the limited sense that the finite pairing
and norm data give it the role of a domain on which Hilbert language might be
extended. It is completion-facing because it is meant to be read inside a
larger object if the door hypotheses are granted. It is not dense by notation,
complete by wish, or continuum-valued by proximity to analytic vocabulary.

One may think of \(\mathcal{S}_{0}\) as the chapter's inventory table. Each
row is a finite record. One column carries its natural-number code. Other
columns carry support, extension behavior, pairing data, and residue data.
There may also be columns reserved for later readings: dense-in, complete-in,
paired-in-the-limit. Those reserved columns are not empty because the chapter
forgot them; they are empty because filling them would require exactly the
completion structure the door has not yet granted. The table format is useful
because it lets the reader see what has a value now and what is only a named
slot for later discharge.

This compact notation also prevents a common ambiguity. The chapter will later
speak about continuum certificate families. Those families require addresses
in the completion: places where geometric and gauge witnesses would sit as
vectors. The skeleton here is not yet such a family of continuum addresses. It
is the countable finite family from which such addresses would have to be
organized. The next section can therefore ask how certificate families should
be named once a completion interface is on the table. It need not reargue why
there is a countable finite family available, and it need not pretend that
availability already provides the continuum addresses.

The separable pre-Hilbert completion skeleton is thus an honest threshold
object. It carries countability, finite extension, compatibility, finite
decision, extensional coding, finite pairing, a zero-detecting norm, and the
finite gauge residue. It names density and completion as obligations still
waiting on the far side of the door. Its strength is exactly this separation.
The finite apparatus arrives with a countable disciplined family in hand; the
completion theory, if admitted, must say how that family becomes dense, how its
limits are housed, and how its finite pairings become continuum pairings. Until
those obligations are discharged, the skeleton remains what the chapter says it
is: a countable pre-Hilbert candidate brought to the threshold, not a completed
Hilbert or Sobolev theorem.

\section{Continuum Certificate Families}

The skeleton of the previous section is a countable finite family brought to
the completion door. It has records, codes, finite pairings, and finite residue
content. This section asks what those records would become if the door supplied
continuum addresses. The answer is a certificate family: not merely a set of
vectors, and not merely a list of finite witnesses, but a disciplined pairing
of finite witness content with a door-side address whose analytic placement is
named rather than built.

A continuum certificate therefore has two components:
\[
  \text{certificate}
    =
    \big(\text{door-side address},\ \text{finite witness}\big).
\]
The second component is the part already earned. It is a finite witness from
the earlier apparatus: a geometric witness, a gauge witness, a generation rule,
or a finite residue receipt. The first component is the part the completion
would have to house. It is called an address rather than a point of an already
constructed space because the chapter is still standing at the threshold. The
address says where the witness would live if the completion hypotheses were
granted. It does not by itself prove that such a completed placement has been
constructed.

This wording matters. A finite witness is content. It says what has been
certified. A continuum address is placement. It says where that content would
sit in a completed geometry. The apparatus has content in hand. It has the
finite geometric record, the finite gauge record, the boundary-generation
relation, and the mixed-pair residue. What it does not yet have is a dense
embedding of those witnesses into a completed Hilbert or Sobolev object. The
certificate notation must therefore keep content and placement separable:
\[
  \text{finite witness} \neq \text{completed vector address}.
\]
The inequality is not a denial that the two should eventually be related. It
is a guard against pretending that the relation has already been supplied.

This separation also keeps the countable skeleton from doing too much work.
The skeleton gives an indexable supply of finite records. A certificate family
adds roles to those records: some records certify geometry, some carry the
generated gauge content, and paired records remember how the gauge side was
generated. What the skeleton still does not add is the continuum map that would
send those records into a completed object with nontrivial positions. Thus the
certificate family is one step richer than the bare skeleton and one step short
of a completed analytic family. It has named witnesses sorted by role; it does
not yet have a completed geometry of their addresses.

The geometric family is the first certificate family to name. Its finite
witnesses come from the geometric side of the earlier chapters: the
energy-and-boundary record, the finite support, the boundary reading, and the
geometric role that organized the certificate sector. These witnesses are not
anonymous labels. They carry the finite structure that made the geometric
sector a certificate sector in the first place. A geometric continuum
certificate would place such a witness at a door-side address:
\[
  G_i = \big(g_i^{\mathrm{addr}},\ w_i^{\mathrm{geom}}\big).
\]
Here \(w_i^{\mathrm{geom}}\) is the finite geometric witness. The address
\(g_i^{\mathrm{addr}}\) is the named completion-side placement. The notation
does not claim that the placement is already nonzero, dense, or complete. It
only gives the future completion theory a precise slot to fill.

The gauge family is named in parallel, but not as an independent decoration.
Its finite witnesses come from the gauge side: the generated gauge program, the
two-channel bookkeeping, the closed linear/vacuum part, and the mixed-pair
residue that survives as charge. A gauge continuum certificate would have the
same two-component form:
\[
  Q_i = \big(q_i^{\mathrm{addr}},\ w_i^{\mathrm{gauge}}\big).
\]
Again, the witness is finite content and the address is completion-side
placement. The section does not say that \(q_i^{\mathrm{addr}}\) already lives
as a rich vector with all analytic relations in place. It says that if the
door is crossed, this is the kind of address the gauge witness must receive.

The two families are related because the earlier apparatus related them. The
gauge witness is not merely adjacent to the geometric witness; it is generated
from the geometric side by the boundary rule. That generation link is finite
content, so it belongs inside the witness side of the certificate. For a paired
record the intended form is
\[
  \begin{aligned}
  (G_i,Q_i)
    &=
      \big((g_i^{\mathrm{addr}},w_i^{\mathrm{geom}}),
          (q_i^{\mathrm{addr}},w_i^{\mathrm{gauge}})\big),\\
  w_i^{\mathrm{gauge}}
    &\text{ is boundary-generated from } w_i^{\mathrm{geom}}.
  \end{aligned}
\]
The relation on the right is the earned finite relation. The addresses on the
left are named positions in the completion interface. This is the correct
division of labor: the generation link is carried forward as a fact; the
continuum placement of the related pair remains an obligation.

The pairing of families should therefore not be read as a product of two
completed vector families. It is a product of two witness families together
with two requested address families. The finite relation says which geometric
witness gives rise to which gauge witness. A completed theory would then have
to place the two addresses in a way that respects that relation. It would have
to say how nearby, orthogonal, trace-compatible, or boundary-related those
addresses are. Those words belong to the completion side. The present section
only ensures that when those words are later used, they will attach to the
right finite witnesses rather than to anonymous vectors.

The placeholder issue can now be stated without embarrassment. The apparatus
can always use the zero address as a formally available address at the door.
That use does not make the certificate trivial, because the finite witness
still carries the content. But it also does not make the zero address the true
continuum placement of the certificate. The zero address is a placeholder used
when no analytic embedding has been granted:
\[
  \begin{array}{c}
  \text{available address now} = 0,\\
  \text{rich certificate placement}
    = \text{named completion upgrade}.
  \end{array}
\]
The witness says what is real on the finite side. The placeholder says what is
not yet real on the completion side. The notation is honest only if both
sentences remain visible.

The placeholder also preserves the distinction between existence of a slot and
occupation of the slot. The door interface can reserve a place for a geometric
address and a gauge address even when the only universally available address is
zero. Reserving the place is useful: it lets later hypotheses state what must
be upgraded. Occupying the place with a nonzero vector is a stronger act. It
would require an embedding, a completed norm, and enough density to make the
certificate family more than a formal placement. The chapter is allowed to
reserve the slot now. It is not allowed to pretend the slot has been occupied
with analytic substance.

This honesty blocks two opposite mistakes. The first mistake would be to throw
away the certificate families because their rich addresses have not yet been
constructed. That would forget the finite witnesses, generation rules, and
residue receipts the apparatus actually earned. The second mistake would be to
promote the placeholder addresses into a continuum geometry. That would invent
angles, distances, nonzero placements, and dense behavior that the door has not
provided. The section takes neither route. It says: the families are real as
finite witness families; their addresses are provisional until the completion
hypotheses supply the analytic placement.

This is why the word ``family'' is still appropriate. The finite witnesses are
not one collapsed object. They remain indexed by their finite records and
their roles. A geometric witness may certify one finite boundary record; a
gauge witness may carry the generated residue for the corresponding record.
The address placeholders do not erase that indexing. They simply fail to add
continuum separation. So the apparatus can speak of two certificate families
without claiming that it has already distributed them across a completed
space. Family structure is earned by the finite ledger and generation rule;
continuum distribution remains named.

The distinction also explains why this section does not yet prove
orthogonality. Orthogonality is a statement about a pairing in the completed
object. The finite apparatus has a sector split and a finite cross-pairing
discipline. It can carry those facts as witness content. But to say that
geometric and gauge certificate vectors are orthogonal in the continuum
requires the continuum pairing and the completed addresses that S01 listed as
door-side obligations. S04 will ask for exactly that interface: orthogonality,
Clifford compatibility, and boundary trace. S03 prepares the families to which
those conditions would apply. It does not silently satisfy the conditions.

The same is true of boundary trace. The finite witness may remember a boundary
reading and the residual charge. It may remember that the gauge witness was
generated from the geometric side. What the present section does not supply is
a Sobolev boundary trace operator on a completed space. The trace is a later
analytic reading of the address, not a property of the witness alone. This
separation keeps the finite boundary record intact while preventing it from
becoming a completed boundary theorem by a change of vocabulary.

The same caution applies to density. A countable family of certificates is not
automatically dense in anything. Countability came from finite coding. Density
would say that completed addresses approximate all the points the relevant
space requires. That sentence needs a space, a metric or norm, and a completed
address map. The present section deliberately stops before that sentence. It
prepares a countable family that could be tested for density once the door
hypotheses are supplied; it does not claim the test has been passed.

One can summarize the section's status in a small ledger:
\[
  \begin{array}{c|c|c}
  \text{fam.} & \text{finite content} & \text{address obligation} \\
  \hline
  G & \text{geometric witness} & \text{geometric placement} \\
  Q & \text{gauge witness/residue} & \text{gauge placement}
  \end{array}
\]
The left and middle columns are available to the finite apparatus. The right
column is the door-facing slot. A genuine completed theory would have to fill
that slot with nontrivial addresses, compatible pairings, and whatever density
or trace behavior the later claims require. The present section earns only the
right to name the slot without confusing it with the witness.

The ledger also shows the order of dependence. First come finite witnesses.
Then comes the relation between witnesses: geometric content generates gauge
content at the boundary. Then come address slots, one for each side of the
relation. Only after those slots are filled can the chapter ask for continuum
relations among the addresses. This order prevents S04 from floating free. Its
orthogonality and trace requirements will not be requirements on a vague
continuum; they will be requirements on the very address slots named here.

The metaphysical anchor is that a certificate is not exhausted by its address.
It has content and placement, and the finite apparatus owns only the content
side for now. The geometric and gauge families carry real witnesses, a real
generation link, and a real finite residue account. Their continuum addresses
are provisional marks in the completion interface. That is enough to let the
chapter proceed to the next question: if such addresses are granted, what must
orthogonality, Clifford compatibility, and boundary trace require of them?

\section{Orthogonality, Clifford Compatibility, and the Boundary Trace}

The previous section named two certificate families. Each certificate carried
finite content on one side and a door-facing address obligation on the other.
This section asks what must be true of those addresses before the finite
boundary-radiation pattern can be read as a continuum statement. The answer has
three parts. The geometric and gauge addresses must be orthogonal in the
completed pairing. The Clifford action must remain compatible with that
orthogonality, so that the mixed scalar obstruction survives the passage to the
door. The generated gauge witness must admit a boundary trace whose value
matches the finite residual it was built to carry. These are not new theorems
proved by the finite ledger. They are the named compatibility obligations that
the next section will place in the antecedent of its conditional result.

The order matters. S03 did not hand us a completed Hilbert space full of rich
vectors. It handed us certificate records:
\[
  \begin{array}{c}
  \text{finite witness content} \\
  \text{door-side address obligation}
  \end{array}
\]
The finite witness content is real. It includes the finite sector split, the
finite Clifford obstruction, and the finite boundary residue produced by the
generation rule. The address obligation is different. It says where the
certificate would have to live once a completion is supplied. The present
section therefore speaks in a double register. It preserves the finite shadows
as genuine inherited content, and it names the analytic conditions that would
make those shadows substantive after completion.

Start with orthogonality. In the finite theory the geometric and gauge sectors
are separated by type. Their cross-pairing vanishes because the finite
apparatus can distinguish the two kinds of witness before any analytic limit is
introduced. This is an earned finite fact. What S04 needs at the door is the
corresponding condition on the addresses:
\[
  \langle g_i^{addr}, q_j^{addr} \rangle_{\mathcal H} = 0
  \qquad\text{whenever the completed pairing is defined.}
\]
The notation is intentionally conditional. It presupposes a completed pairing
\(\langle -, - \rangle_{\mathcal H}\) and address placements for the
geometric and gauge certificates. The finite construction supplies neither a
completed inner product nor a dense embedding theorem here. It supplies the
finite split whose completed counterpart would have to satisfy the displayed
orthogonality equation.

The zero placeholder clarifies the grammar but not the substance. If both
addresses are interpreted as the distinguished zero address, then the displayed
equation has a formal shadow: zero pairs to zero. That shadow is useful because
it shows that the address slot has the right type of shape and that the
finite-to-door bookkeeping is coherent. It is not the continuum orthogonality
claim the chapter ultimately cares about. A nontrivial orthogonality claim
would require nonzero or otherwise rich placements of the geometric and gauge
families, a completed pairing on those placements, and a proof that the pairing
vanishes on the intended cross-terms. The finite split tells us exactly what
the condition should be. It does not discharge that condition.

There is a further reason to state the orthogonality obligation at the level of
addresses rather than at the level of witnesses. A finite witness can remember
which sector it came from, but an analytic vector can forget its origin unless
the placement map preserves it. The obligation therefore has two pieces. First,
there must be a placement map that sends geometric witnesses to geometric
addresses and gauge witnesses to gauge addresses without collapsing the
distinction. Second, the completed pairing must respect that distinction by
annihilating the cross-pairings that the finite split already knows how to
annihilate. Written schematically, the completed theory must make the square
\[
  \begin{array}{ccc}
  \text{finite sector split} & \longrightarrow & \text{finite zero cross-term} \\
  \downarrow & & \downarrow \\
  \text{address placement} & \longrightarrow & \text{completed zero cross-term}
  \end{array}
\]
commute in the only sense relevant here: the zero cross-term seen finitely must
remain zero after placement. This is not a density assertion. It is a
preservation assertion.

This distinction prevents two common mistakes. The first mistake is to treat
the finite sector split as if it automatically generated orthogonal vectors in
a completed space. The second is to dismiss the finite split because it does
not yet do so. The right reading sits between them. The finite split is the
reason the orthogonality obligation has a canonical form. The completion
hypotheses are what would make that form analytically true. Thus the
orthogonality requirement is neither an empty slogan nor a completed result; it
is an exact address-level condition carried forward from finite content.

The second obligation is Clifford compatibility. Earlier finite chapters used
the Clifford relation to turn orthogonality into an obstruction. In the finite
setting, the mixed scalar part is governed by the same pairing that separates
the sectors. When the pairing cross-term is zero, the interior mixed scalar
certificate has nowhere to live. At the door, this becomes a compatibility
condition between three pieces of structure:
\[
  \begin{array}{c}
  \text{completed pairing} \\
  \text{Clifford action} \\
  \text{mixed scalar certificate rule.}
  \end{array}
\]
The condition says that the Clifford action must read the completed pairing in
the same way the finite obstruction read the finite pairing. In particular,
\[
  \langle g_i^{addr},q_j^{addr}\rangle_{\mathcal H}=0
\]
should prevent the completed Clifford interpretation from reintroducing an
interior mixed scalar certificate by another route.

One compact way to state the obligation is:
\[
  \langle x,y\rangle_{\mathcal H}=0
  \quad\Longrightarrow\quad
  \operatorname{mix}_{Cl}(x,y)=0,
\]
where \(\operatorname{mix}_{Cl}\) names the scalar mixed component read through
the chosen Clifford-compatible structure. This formula should be read as a
requirement on a future completed structure, not as an assertion that the
structure has already been constructed. The finite theory supplies the pattern:
orthogonality kills the scalar obstruction. The completion must supply the
objects on which that pattern can act: address placements, a completed pairing,
a Clifford action, and a proof that the action respects the pairing in the
needed way.

The placeholder again gives only a formal shadow. At the zero address the mixed
term may vanish for the same reason every bilinear or scalar term vanishes at
zero. That is not the obstruction theorem. The obstruction theorem is about
nontrivial geometric and gauge content retaining their separation after the
door is crossed. S04 therefore refuses to say that the Clifford compatibility
has been proved. It says that the finite Clifford obstruction determines the
compatibility condition the completed theory must satisfy. The condition is
sharp enough to be useful: if a proposed completion allows the mixed scalar
term to reappear despite the orthogonality condition, then it has failed to
preserve the finite obstruction.

The compatibility requirement also explains why orthogonality alone is not
enough. A completed pairing could make the geometric and gauge addresses
orthogonal while a separately chosen Clifford action failed to respect that
pairing. Such a structure would have the right inner-product sentence and the
wrong algebraic behavior. The finite result tied the two together: the scalar
part of the Clifford relation was controlled by the same pairing that expressed
sector separation. The door must keep that tie. Thus the completed Clifford
action is not merely an extra decoration placed on the completed space. It is a
test of whether the completed space has remembered the finite obstruction in
the right algebraic language.

For this reason, the mixed scalar certificate is the best diagnostic object. If
it appears in the interior, the completion has lost the finite story. If it is
blocked only because all addresses were replaced by zero, the completion has
not yet gained the substantive story. The desired situation is sharper: rich
geometric and gauge addresses are present, the completed pairing makes the
intended cross-term vanish, and the Clifford action converts that vanishing
into the absence of an interior scalar mixed certificate. The section names
that situation because S05 will need it as a hypothesis.

The third obligation is the boundary trace. The finite boundary pattern has a
generated gauge witness whose residual is read at the boundary.
In the finite language this residual is terminal: it is what remains after the
interior mixed certificate is obstructed. At the door the same fact must be
expressed through a trace:
\[
  \operatorname{Tr}(q_i^{addr}) =
  \operatorname{res}(w_i^{gauge})
  \qquad\text{for generated gauge certificates.}
\]
Here \(\operatorname{res}(w_i^{gauge})\) belongs to the finite witness content.
The trace \(\operatorname{Tr}\) belongs to the analytic side. The equation
therefore crosses the same boundary as the rest of the chapter: finite residue
on one side, completed boundary functional on the other.

This trace requirement is deliberately modest. It does not assert a Sobolev
space has been completed. It does not assert that a trace operator exists for
all relevant elements. It does not identify the finite residual with a
physical flux by fiat. It says what must be true if the generated gauge witness
is to be read as boundary radiation in a completed setting: the completed
address of that witness must lie in the domain of a trace, and the trace value
must preserve the finite residual assigned by the generation rule. The finite
construction owns the residual and the generation link. The analytic
completion must own the domain, the trace map, and the preservation statement.

The trace obligation has its own preservation square. The finite construction
starts with a generated gauge witness and a residual value. The analytic side
must start with the address of that generated witness and a trace map. The
square says that tracing the address must recover the same residual that the
finite construction assigned:
\[
  \begin{array}{ccc}
  w_i^{gauge} & \longmapsto & \operatorname{res}(w_i^{gauge}) \\
  \downarrow & & \uparrow \\
  q_i^{addr} & \longmapsto & \operatorname{Tr}(q_i^{addr})
  \end{array}
\]
The arrows here are explanatory rather than new machinery. They mark the
preservation demand. The finite side does not construct the lower horizontal
arrow, and the analytic side may not rewrite the upper value. Agreement is the
point: the boundary reading must preserve the residual, not replace it with a
new quantity that merely resembles it.

This is also where the word ``boundary'' earns its restraint. The finite
residual was already boundary-facing, but a boundary-facing finite record is
not yet a trace theorem. A trace theorem would require a domain, a boundary
operator, and a statement that the generated representative lies in that
domain. Until those are supplied, the residual remains finite content awaiting
an analytic reader. The chapter can say exactly what the reader must do; it
cannot claim the reader exists just because the finite record is ready for it.

The generation link is the thread that ties the three obligations together. A
geometric witness does not merely sit beside a gauge witness. Earlier chapters
made the gauge witness arise from the geometric side as the boundary-facing
residue of the mixed obstruction. S03 preserved this relation at the level of
finite witness content. S04 now asks the completed addresses to preserve it as
well. Orthogonality separates the two address families. Clifford compatibility
explains why the interior mixed scalar certificate remains obstructed. The
trace condition says where the generated residual is allowed to appear: at the
boundary, as the value of the trace on the generated gauge representative.

This gives the section a useful ledger:
\[
  \begin{array}{c|c|c}
  \text{finite shadow} & \text{door obligation} & \text{analytic content} \\
  \hline
  \text{sector split} & \text{orthogonality} & \text{completed pairing} \\
  \text{Clifford obstruction} & \text{compatibility} & \text{Clifford action} \\
  \text{boundary residue} & \text{trace agreement} & \text{trace operator}
  \end{array}
\]
The first column is inherited. The second column is what the chapter now names.
The third column is what a nontrivial completion must provide. This table is
also a guard against overclaiming: none of the items in the third column is
created merely by naming it. The table records obligations, not completed
analytic infrastructure.

One can now see why the three obligations must be stated together rather than
as independent wishes. Orthogonality without Clifford compatibility would give
a silent pairing but no obstruction. Clifford compatibility without a trace
would block the interior mixed scalar term but leave the residual with no
boundary reading. A trace without orthogonality and Clifford compatibility
would read a residual whose interior origin had not been controlled. The finite
boundary-radiation result had all three pieces at once: separation,
obstruction, and boundary residue. The door-facing version must keep the same
simultaneity.

The simultaneity is especially important for the generated gauge witness. The
gauge witness is not an arbitrary element selected after the fact. It is the
one produced by the finite generation relation from the geometric side. The
completed theory therefore has to preserve not only a pair of isolated
addresses but also the fact that the gauge address is the address of the
generated witness. Otherwise a completion could satisfy a trace equation for
some gauge-like element while losing the specific finite source of that
element. The chapter refuses that shortcut. The boundary trace must apply to
the generated representative, and the generated representative must still be
linked to the finite obstruction story.

The zero placeholder belongs in the ledger only as a diagnostic. It can show
that the formulas are well formed at the door. It can confirm that the record
has a place for a geometric address, a gauge address, a Clifford reading, and a
trace reading. It can even make certain equations formally true in the most
degenerate case. But a degenerate witness is not the theorem. The theorem
would require rich addresses carrying the finite content through a completion
without losing the sector split, the Clifford obstruction, or the generated
residual. The placeholder tells us that the slot is coherent. It does not fill
the slot with the continuum object.

There is a productive use of the placeholder nevertheless. It allows the
chapter to separate type discipline from analytic achievement. Type discipline
asks whether the records have the right fields, whether the formulas mention
the right objects, and whether the equations are shaped so that a future
completion could interpret them. Analytic achievement asks whether nontrivial
placements, pairings, actions, and traces have actually been supplied. The
placeholder can support the first question. It cannot settle the second. This
separation is what lets S04 be precise without being inflated.

This is why the section keeps returning to preservation. A completion that
forgets the finite split is too coarse. A completion that supplies a pairing
but lets the Clifford action ignore that pairing is incompatible. A completion
that carries gauge addresses but provides no trace, or a trace that fails to
read the generated residual, breaks the boundary-radiation interpretation. The
finite construction gives a compact test suite for any proposed analytic
upgrade: preserve separation, preserve the obstruction, and preserve the
boundary reading. Failure of any one of the three means the conditional theorem
in the next section cannot be invoked.

The result is a clean division of labor. The finite side contributes witnesses,
sector distinctions, the mixed obstruction, and the residual bookkeeping. The
completion side must contribute a completed inner product, nontrivial address
placements, a Clifford action compatible with that inner product, a trace
operator, and preservation of the finite generation link. The chapter does not
pretend those two sides are the same achievement. It uses the first to state
the exact hypotheses required of the second.

Thus S04 is not a proof of continuum orthogonality, continuum Clifford algebra,
or Sobolev trace. It is the place where those three requirements become
explicit and coordinated. S03 gave the certificates addresses. S04 tells those
addresses what they must satisfy. S05 can now state the generalized
boundary-radiation claim honestly: if the address placements are orthogonal, if
the Clifford action is compatible, if the generated gauge witness has the
right trace, and if the remaining completion hypotheses are supplied, then the
finite boundary-radiation pattern has a legitimate continuum reading.

\section{Generalized Boundary Radiation as a Conditional Theorem}

The chapter now reaches the form of the frontier theorem. It is not an
unconditional theorem of analysis, and it is not a proof that the finite
construction has already built the completed boundary space. It is a conditional
theorem with its conditions written on the face of the statement. The preceding
sections have not been leading to a sleight of hand in which finite data suddenly
become a continuum theorem. They have been preparing an implication: if the
completion supplies the placements, pairings, actions, traces, preservation
facts, and convergence facts that the finite apparatus has isolated, then the
boundary-radiation pattern produced in the finite theory has a legitimate
continuum reading.

That distinction is the point of the section. The finite work does something
real. It supplies the witnesses, the sector split, the obstruction pattern, the
generated gauge representative, and the exact shape of the analytic hypotheses
that must be added at the door. It identifies which certificates are meant to be
placed in the completion, which pairings are meant to become orthogonal, which
action is meant to respect the Clifford structure, which generated
representative is meant to carry the trace, and which terminal residue is meant
to survive as boundary flux. What it does not supply is the analytic
infrastructure itself. It does not prove density, completeness, nonzero
orthogonal embeddings, a continuum Clifford algebra, or a trace theorem. It
separates the finite contribution from the analytic contribution and then states
the theorem that follows when both are present.

The antecedent has six parts. First, the finite certificates must have address
placements in the chosen completion setting. The proton-side and electron-side
finite witnesses cannot remain merely formal labels if the theorem is to be read
as a statement about a completed boundary space; they must be placed as objects
of that space. Second, the completed pairing must make the two placed sectors
orthogonal in the required sense. This is the continuum version of the separated
sector pattern produced earlier, and it is not discharged by saying that two
finite names are different. Third, the action on the completed objects must be
compatible with the Clifford structure. The finite construction names the
sign-bearing action and the way it interacts with the certificates, but the
completed action must be supplied as part of the analytic setting.

Fourth, the completion must provide a boundary trace operator in the relevant
class. The generalized radiation reading needs an operator that can take the
generated representative to the boundary expression that the finite theory has
singled out. Fifth, the link between finite generation and the completed
representative must be preserved. The residue in the completed setting is not an
arbitrary boundary object; it is the continuation of the same generated gauge
representative that arose from the finite obstruction. Sixth, the residual
sequence must satisfy the named convergence condition. In the present chapter
that condition is carried as a Cauchy condition in the completion setting. The
finite apparatus can say exactly which residual sequence and which terminal
residue are at issue, but the completed norm and its convergence facts belong to
the analytic side of the door.

With those hypotheses named, the conditional theorem can be stated. Suppose that
the certificates have completion addresses, that the completed pairing separates
their sectors orthogonally, that the completed action is Clifford-compatible,
that the boundary trace is defined and agrees with the terminal residual on the
generated representative, that finite generation is preserved by the completed
representative, and that the residual sequence is Cauchy in the named completion
setting. Then the generalized boundary-radiation conclusion follows:
\[
  \left.
  \begin{array}{l}
    \text{completion addresses} \\
    \text{orthogonal pairing} \\
    \text{compatible action} \\
    \text{boundary trace} \\
    \text{preserved generation} \\
    \text{residual Cauchy}
  \end{array}
  \right\}
  \;\Longrightarrow\;
  \begin{array}{l}
    \text{no interior mixed certificate} \\
    \text{flux} = \text{terminal residual} \\
    \text{residual convergence}
  \end{array}
\]
The first line of the consequent says that the mixed scalar certificate still
does not appear in the interior. The separation of sectors has not been erased by
passing through the door; under the completed orthogonality hypothesis, the
obstruction to an interior mixed scalar remains in force. The second line says
that the boundary flux is the terminal residual. This is the continuum reading of
the finite radiation pattern: the generated representative does not become
interior matter, but appears at the boundary as the residue singled out by the
trace. The third line says that this residual has the convergence status required
in the named completion. The theorem therefore does not merely repeat the finite
calculation. It says what the finite calculation becomes when the analytic
setting validates the placements, operations, and limiting process that the
calculation asks for.

The theorem is conditional because each hypothesis carries real weight. Address
placement is not cosmetic: without it, the finite certificates
are not objects of the completed setting. Orthogonality is not a slogan: without
the completed pairing, there is no continuum meaning to the sector separation.
Clifford compatibility is not inherited by vocabulary: without a compatible
completed action, the sign-bearing structure of the finite theory has no
licensed analytic extension. Trace compatibility is not automatic: without a
trace operator, the boundary expression is only a formal echo of the finite
residue. Preservation of finite generation is not optional: without it, the
boundary object could have lost the very origin that makes it radiation rather
than arbitrary boundary data. Residual convergence is not decoration: without
it, there is no completed limiting statement to accompany the terminal residue.
The antecedent is the theorem's analytic content.

This is also why the result does not claim more than it has earned. It does not
say that the finite apparatus has produced a Hilbert completion. It does not say
that a Sobolev trace theorem has been proved. It does not say that nonzero
continuum representatives have been constructed or that the desired completed
Clifford action exists. It says that if such a completion is supplied with the
listed compatibilities, then the finite boundary-radiation mechanism transports
to that setting in the precise form displayed above. The price of the theorem is
visible. The benefit of paying that price is visible as well: the finite
obstruction pattern becomes a continuum boundary statement without pretending
that the analytic bridge has already been crossed.

The zero placement used inside the finite formal development should be read in
that limited spirit. It is a formal sanity check and a degenerate witness for
the shape of the implication. At the zero placement, the address data,
orthogonality, compatibility, trace agreement, and residual condition can be
recorded without adding analytic substance. This verifies that the conditional
surface is coherent: the hypotheses fit together, and the consequent has the
intended logical form. It is not the rich theorem one ultimately wants. It is
not a nonzero continuum radiation theorem, not a density result, and not an
analytic construction of the completed boundary theory. It is the finite
apparatus showing that the implication is well-formed and that no hidden
contradiction has been introduced in the act of naming the door.

The substantive continuum theorem would begin where the zero sanity check stops.
It would supply nontrivial address placements, prove the completed pairing
orthogonal for those placements, construct the compatible completed action,
establish the trace operator, show that the generated representative remains the
right representative after completion, and verify residual convergence in the
chosen norm. That work is deliberately outside the present finite theory. The
chapter refuses to smuggle it in under a suggestive name. Instead, it records the
exact analytic tasks that would have to be solved and records the exact theorem
that would then be available.

This makes the finite contribution sharper, not weaker. The finite theory does
not merely gesture toward analysis; it hands analysis a specified boundary
problem. The data are not anonymous. The proton-side and electron-side witnesses
are the ones already separated by the obstruction. The gauge representative is
the one generated by that obstruction. The terminal residual is the same residue
that carried the radiation reading. The requested trace is the trace of that
representative, not of an unrelated boundary field. The requested convergence is
the convergence of that residual pattern, not a generic compactness assertion.
Thus the analytic continuation, if supplied, is constrained by the finite
construction from which it starts. The conditional theorem preserves that
constraint in its statement: every hypothesis points back to a finite feature,
and every conclusion says what that feature becomes after completion.

Seen this way, the finite apparatus has not failed at the continuum door. It has
reached the door with enough structure to say what the door is. Earlier chapters
drove the invariant to a signed charge, produced the obstruction that separates
the sectors, identified the generated gauge residue, and interpreted that
residue as boundary radiation. The present chapter added the discipline needed
at the frontier: address placement, completed orthogonality, Clifford
compatibility, trace compatibility, preservation of finite generation, and
residual convergence. These are not afterthoughts. They are the checklist that
turns a finite radiation pattern into a conditional theorem about a completed
boundary setting.

The closing lesson of the numbered sections is therefore methodological as much
as mathematical. A finite theory can carry a continuum-facing result only by
making the missing analytic content explicit. If it hides the hypotheses, it
overclaims. If it refuses to state the implication, it underuses the structure it
has built. The right form is the conditional: finite witnesses and finite
obstructions on one side, analytic completion hypotheses on the other, and a
precise boundary-radiation conclusion following from their conjunction. That is
what the door names. It names neither a completed crossing nor an empty wish, but
the exact interface between the finite mechanism already constructed and the
analytic theory that would let that mechanism speak in the continuum.

Thus Chapter 16 closes its numbered work by fixing the status of the frontier
claim. The apparatus has not proved an unconditional generalized
boundary-radiation theorem. It has proved the conditional shape that any such
theorem must respect, identified the finite data that feed it, and named the
analytic obligations that remain. Once those obligations are supplied, the
boundary flux is the terminal residual, the mixed scalar certificate is not an
interior object, and the residual has its convergence reading in the completion.
Until then, the door remains honestly closed: visible, specified, and ready for
the analytic work it requires.

\section*{Bridge: One Convention Left to Anchor}

Chapter 16 has finished the continuum-facing work it was allowed to do. It has
not crossed the Sobolev/Hilbert door, and it has not pretended that a finite
ledger secretly contains a completed analytic theory. Instead it has brought
the apparatus to the threshold with its credentials in order: a finite norm and
pairing, a countable skeleton, certificate families with finite witness
content, compatibility obligations for orthogonality, Clifford action, and
trace, and finally a conditional boundary-radiation statement. The chapter's
achievement is the shape of the door. It names the price of crossing and
refuses to claim that the price has already been paid.

That refusal matters because it keeps the book's central habit intact. At each
frontier, the apparatus has distinguished what it built from what it merely
needed. It did not call a finite code a continuum. It did not call a zero
placeholder a nontrivial embedding. It did not call a carried trace condition a
trace theorem. It did not call a conditional theorem an unconditional theorem.
The result is narrower than a triumphant analytic completion, but cleaner: if a
completion supplies the named structures, the finite boundary-radiation pattern
has a legitimate reading there; until then, the statement remains conditional.
The door is visible, specified, and honestly closed.

With that door named, the book can no longer postpone the one convention it has
carried from the moment the pair appeared. The apparatus produced a signed
charge. Its magnitude was not the problem. The finite construction could say
that the mixed pair leaves one unit of charge, and the later chapters could
carry that unit through boundary generation and radiation. But a signed unit
has an orientation. Read one way, the charge is the electron-side value; read in
the opposite orientation, the mirror-side value appears. The apparatus has used
the electron orientation throughout the book. What it has not yet done is say
what kind of act that choice is.

The point is not that the sign was forgotten. It has been present all along,
quietly doing work. The signed charge let the apparatus distinguish the two
members of the pair. It let the boundary story speak as electron rather than as
its mirror. It let later sections keep a coherent physical reading while the
formal machinery remained finite. But coherence is not the same as derivation.
The fact that a convention is stable across the argument does not by itself
make the convention forced. Chapter 16 has just taught the book to be careful
about this distinction. If analytic completion must be named as an obligation
rather than smuggled into a theorem, then orientation must be treated with the
same honesty.

The final chapter therefore does not reopen the completion door. It turns from
the question of crossing a boundary to the question of choosing an orientation.
It asks what remains after magnitude has been separated from sign. A magnitude
can be produced as the size of the finite residue. A sign requires an ordered
reading of the two possible directions. The apparatus can display both
orientations, but it still has to account for why one is carried as the
electron-side convention and the other as its opposite. That account cannot be a
new theorem smuggled in at the end, and it cannot be a shrug. It has to say how
a convention can be anchored without being derived.

This is the last turn of the book's method. Earlier chapters showed that an
instrument does not need an omniscient outside view in order to make a stable
record. It needs bounded repeatability, a way to distinguish signal from
discard, and a disciplined account of what its own procedure has fixed. The
orientation problem is the same kind of problem at the sign level. A finite
apparatus may be unable to derive an absolute orientation from nothing, but it
can measure which orientation dominates inside its own arrangement once the
noise floor has been named. The sign convention is then not metaphysical
legislation. It is an anchored reading: a deliberate settlement of which side
of a symmetric pair the apparatus will call electron-side.

Chapter 17 will take up that settlement. It will distinguish an orientation
from its mirror, separate the magnitude of the charge from the sign by which
the magnitude is read, and ask how dominance past a noise floor can anchor a
choice without turning that choice into a derivation. The chapter must be
careful for the same reason Chapter 16 was careful. To anchor a convention is
not to prove that the opposite convention is impossible. It is to show why this
apparatus, with this record and this orientation-sensitive measurement, can
carry one convention responsibly.

So the bridge out of the completion door is narrow. The book has named the
analytic price of a continuum reading and stopped short of paying it. Now it
names the remaining classical price of its own signed charge: the orientation
must be faced directly. The apparatus has a charge with a magnitude, and it has
used one sign convention to read that charge. The final question is what it
means for a measurement to anchor that convention rather than derive it from
nowhere.

\section*{Coda: The House at Escrow}

A house under contract but not yet closed gives Chapter 16 its cleanest image.
The buyer and seller have done something real. They have not merely talked
about a possible sale. They have signed a contract. The price is written down.
Earnest money has gone into escrow. The seller has taken the house off the
market. The buyer has shown enough seriousness that the deal has a shape, a
date, a file number, and a closing table waiting somewhere ahead. Everyone
involved can point to an actual commitment.

And yet nobody owns the house differently at that moment. The buyer does not
walk out of the signing with keys. The seller does not hand over the deed
because the word ``contract'' has appeared on paper. Between signing and closing
stand conditions precedent. The inspection must come back clean or be settled.
The buyer's financing must be approved in fact, not merely hoped for. The title
search must show that the seller can convey what the seller has promised to
convey. The appraisal may have to support the loan. The closing funds must
arrive. Each condition is ordinary, concrete, and named. Each is part of the
deal because the deal is serious enough to say what could still stop it.

That is the structure of the chapter. The finite apparatus has signed a real
contract with the continuum door. It brings genuine credentials: a
zero-detecting norm, a finite pairing, a countable skeleton, finite certificate
families, a generated gauge residue, and a conditional boundary-radiation
claim. These are not wishes. They are the earnest money of the theory, the
signed purchase agreement, the seller's promise to hold the house for this
buyer. The finite work has made a serious commitment to a particular analytic
future. It has not bought every possible house; it has named this house, this
price, this route to closing.

But the analytic closing has not happened. The completed placement of the
certificates is still an inspection item. The completed pairing is still a
title question: does the structure actually convey the separation that the
finite record promises? The compatible action is still a financing condition:
can the sign-bearing mechanism survive the completed setting rather than only
the finite application? The boundary trace is still a required report: can the
generated representative be read at the boundary in the promised way? Residual
convergence is still money that must arrive at closing, not a letter saying
that it probably will. The chapter lists these conditions because the contract
would be dishonest without them.

Escrow is the right image for the placeholders. Earnest money is real money,
but it is not the purchase price delivered to the seller. A pre-approval letter
is a real credential, but it is not funded financing. A title commitment is a
real report, but it is not yet the deed recorded in the buyer's name. These
documents and deposits are not fake. They are held precisely because the deal
is serious. They mark what has been committed and what remains to be
discharged. The zero placement and completion tokens in the chapter play that
role. They keep the shape of the theorem visible while refusing to pretend that
the analytic substance has already arrived.

A bad reading of escrow would say: because the money is in escrow, the house is
already bought. A bad reading of Chapter 16 would say: because the finite
apparatus has named the completion and displayed placeholders, the continuum
theorem has already been proved. Both readings confuse commitment with
closing. The earnest money makes the agreement binding; it does not satisfy the
inspection, fund the loan, clear the title, or record the deed. The finite
credentials make the frontier theorem precise; they do not supply density,
completion, nonzero orthogonality, a trace theorem, or a completed action.

The conditional contract is therefore not a defect in the sale. It is the sale
written honestly. It says: if the inspection passes, if financing clears, if
title is clean, if funds arrive, then title transfers. The word ``if'' is not a
retreat from seriousness. It is the legal form of seriousness before closing.
It protects the buyer from pretending that a defective title is clean. It
protects the seller from pretending that promised funds have arrived. It tells
every party what must happen, and it tells them what follows once it does.

The chapter's theorem has that same legal grammar. If the completion supplies
addresses, if the pairing is completed in the required way, if the action is
compatible, if the boundary trace exists and agrees with the generated
representative, if finite generation is preserved, and if the residual has its
convergence reading, then the boundary-radiation pattern has a legitimate
continuum interpretation. The mixed scalar certificate is not an interior
object. The boundary flux is the terminal residual. The residual belongs to the
named completion story. The theorem is not weakened by the antecedent. The
antecedent is the closing package.

One can feel the temptation to rush the signing table. The house is beautiful;
the buyer wants to move in; the seller wants the transaction finished; the
contract is already thick with signatures and initials. It would be satisfying
to take a photograph on the front porch and call the matter done. But a careful
transaction does not work that way. It honors the difference between under
contract and closed. It lets the inspectors inspect, the lender lend, the title
company search, the escrow officer hold, and the closing agent record. That
patience is not bureaucracy for its own sake. It is the discipline that makes
the final transfer real.

Chapter 16 asks for the same patience. The finite apparatus has enough to stand
at the threshold and say what would happen if the analytic conditions were
met. That is a strong result. It is stronger than a vague aspiration toward a
continuum theory, because every condition has been named. It is stronger than a
mere metaphor, because the finite witnesses and residues are actually in hand.
But it remains weaker than closing. It does not hand itself the keys. It does
not move the furniture into a completed Hilbert space. It waits at escrow with
the contract signed and the conditions written.

This waiting is not empty. Escrow is an active form of waiting. Documents are
collected. Reports are reviewed. Funds are verified. Defects are either cured
or allowed to stop the sale. The buyer's seriousness is measured by staying in
the process without pretending the process has ended. Likewise, the apparatus's
seriousness is measured by holding its finite construction in view while
leaving the analytic obligations visible. It does not withdraw the offer, and
it does not forge the deed. It keeps the conditional alive.

The house image also explains why the chapter's placeholders are useful even
when they are not the closing itself. Escrow records who owes what. It holds
the earnest money so that the promise is not weightless. It keeps the contract
from dissolving into conversation. The completion token and zero placement do
the same. They mark the slots where analytic work must enter. They prevent the
finite theorem from losing its shape while the missing conditions are named.
They are not the finished transfer, but without them the closing package would
not even be organized.

There is another virtue in the metaphor: it separates confidence from
possession. The buyer may be confident after the inspection. The lender may be
nearly ready. The title company may have found only routine matters. Everyone
may expect the closing to happen next week. Still, expectation is not title.
The signatures are gathered because the parties expect a transfer, but the
transfer waits for the required acts. Chapter 16 lives in that same interval.
It has enough structure to make the conditional theorem worth stating. It has
enough finite evidence to say why these are the right conditions rather than an
arbitrary wish list. But expectation, plausibility, and formal placeholders are
not analytic possession. The keys remain on the table until the conditions are
met.

This is why the coda should not be read as a retreat from the chapter's
mathematics. A contract that names conditions precedent is not less serious
than one that leaves them vague. It is more serious. It tells the parties where
to look, what to repair, and what would count as completion. Likewise, the
chapter's conditional theorem is not an apology for failing to speak about the
continuum. It is the exact form in which the finite apparatus can speak about
the continuum without overclaiming. It says: here is the finite record; here is
the analytic closing package; here is what title transfer would mean.

And if a condition fails, the contract has still done useful work. A failed
inspection does not make the signed contract meaningless; it reveals the defect
that must be cured or the reason the sale cannot close. A title problem does
not erase the earnest money; it explains why title cannot yet pass. The named
conditions protect the parties from confusion. In the same way, if the analytic
completion cannot supply one of the named obligations, the chapter has not
spoken falsely. It has located the failure. It has said which part of the
closing package is missing and why the unconditional reading cannot yet be
claimed.

So the honest sentence at the end of Chapter 16 is not ``the house is ours.''
It is: the house is under contract, the earnest money is in escrow, the
conditions precedent are listed, and if they are satisfied, title passes. That
sentence is conditional, but it is not timid. It says exactly what has been
done and exactly what remains. It respects the signed agreement without
confusing it with possession. It respects the future closing without pretending
that future work has already occurred.

That is the door. The finite apparatus has signed the contract. The completion
conditions are in the file. The placeholders sit in escrow. The conditional
theorem says what closing would buy. Until the inspection, financing, title,
and funds have their analytic counterparts, the keys stay on the table. The
deal is real. The house is not yet closed. Chapter 16 ends there, because that
is the only place an honest threshold chapter can end.
